\documentclass[10pt,a4paper]{article}
\usepackage[T1]{fontenc}\usepackage[utf8]{inputenc}\usepackage[english,italian]{babel}
\usepackage{lmodern,microtype}\usepackage[a4paper,margin=1.85cm,headheight=15pt]{geometry}
\usepackage{enumitem,tabularx,array,longtable,booktabs,xcolor,hyperref,xurl,fancyhdr,titlesec,framed}
\definecolor{ddtadark}{RGB}{28,38,52}\definecolor{ddtablue}{RGB}{42,88,135}\definecolor{ddtagray}{RGB}{244,246,248}\definecolor{ddtagreen}{RGB}{48,111,78}\definecolor{ddtaorange}{RGB}{148,91,25}
\hypersetup{colorlinks=true,linkcolor=ddtablue,urlcolor=ddtablue,pdftitle={SRC-0048 ISO IEC IEEE 42010 2011 BA0-R regression}}
\pagestyle{fancy}\fancyhf{}\lhead{DDTA - BA0-R10}\rhead{REVIEW-ONLY}\cfoot{\thepage}
\titleformat{\section}{\large\bfseries\color{ddtadark}}{}{0pt}{}\setlist[itemize]{leftmargin=1.45em,itemsep=.18em,topsep=.2em}
\setlength{\parindent}{0pt}\setlength{\parskip}{.32em}\setlength{\emergencystretch}{2em}
\newenvironment{factbox}[1]{\def\FrameCommand{\fcolorbox{ddtablue}{ddtagray}}\MakeFramed{\advance\hsize-2\width\FrameRestore}\noindent\textbf{#1}\par\smallskip}{\endMakeFramed}
\newenvironment{challenge}[1]{\def\FrameCommand{\fcolorbox{ddtaorange}{ddtagray}}\MakeFramed{\advance\hsize-2\width\FrameRestore}\noindent\textbf{#1}\par\smallskip}{\endMakeFramed}
\begin{document}
\begin{center}{\LARGE\bfseries DDTA - Scheda di lettura compilata}\par
{\large SRC-0048 - ISO/IEC/IEEE 42010:2011 full-text regression}\par
{\small REVIEW-ONLY: historical full-text source; 2022 edition remains separate ACCESS-LIMITED SRC-0043.}\par\end{center}
\begin{factbox}{Identita bibliografica e accesso}
\begin{tabularx}{\linewidth}{>{\bfseries}p{3.6cm}X}
Source ID & SRC-0048\\
Titolo & ISO/IEC/IEEE 42010:2011 - Systems and software engineering -- Architecture description\\
Ente & ISO; IEC; IEEE\\
Edizione & First edition, 2011-12-01\\
Accesso letto & user-supplied full-text research copy\\
SHA-256 & \texttt{0b53b9400cd590f9a89d767333ff6d80f2296892a76ab8e08225eae035edb64f}\\
Repository storage & FORBIDDEN; source PDF excluded from drop-in\\
Verification state & verified\_user\_supplied\_full\_text\\
\end{tabularx}
\end{factbox}
\section{1. Research problem and contribution}
The standard governs architecture descriptions, viewpoints, frameworks and ADLs. It sharply distinguishes the abstract architecture of a system from the work product used to express it. It does not define what a system is and does not prescribe one method, notation or medium.
\section{2. Starting artifacts and representation}
The conceptual model separates system-of-interest, architecture, architecture description, stakeholder, concern, viewpoint, view, architecture model, model kind, correspondence, decision and rationale. Annex B states that a model-kind metamodel can define entities, attributes, relationships and constraints. This supports the BA0-R extraction lens but does not determine the DDTA BAE classes.
\section{3. Traceability, change, automation and human role}
Correspondences explicitly relate architecture-description elements and can be governed by rules. They can express consistency, traceability, composition, refinement, dependency and model transformation, including links to requirements. Key architecture decisions can carry identity, affected elements, concerns, owner, rationale, constraints, alternatives, consequences, timestamps and source citations. Architecture descriptions can aggregate several life-cycle work products and serve as input to simulation, generation and analysis.
\section{4. Main DDTA challenge}
\begin{challenge}{Earlier projection claim is too strong as prior-art claim}
Annex A describes both a synthetic approach, where separately constructed views are integrated through correspondences, and a projective approach, where views are extracted from an underlying repository. ISO 42010 supports either. DDTA may still choose a canonical projective representation, but BA0/BA4 must justify it as a DDTA design decision rather than treat it as mandatory architecture semantics.
\end{challenge}
\section{Hypothesis effects}
\begin{tabularx}{\linewidth}{>{\bfseries}p{4.5cm}X}
Architecture != description & STRONGLY SUPPORTED.\\
View != viewpoint & STRONGLY SUPPORTED.\\
One canonical repository mandatory & REJECTED as an ISO prior-art claim.\\
Projective DDTA views & CANDIDATE design choice, not external invariant.\\
Partial diagram == architecture view & REJECTED if ISO terminology is adopted; partial content fits architecture model/projection better.\\
Correspondence semantics & STRONGLY EVIDENCED prior art.\\
Decision/rationale & STRONGLY EVIDENCED prior art.\\
ISO 42010 supplies BAE ontology & REJECTED; it explicitly does not define system nature.\\
\end{tabularx}
\section{Citation-ready evidence}
\begin{longtable}{p{1.0cm}p{2.1cm}p{4.0cm}p{7.4cm}}
\toprule PDF & Ruolo & Sezione & DDTA interpretation\\\midrule\endhead
8 & foundation & Terms & Architecture, AD, view, viewpoint and model kind are distinct.\\
10 & challenge & 4.2.1 & System nature is intentionally not defined; environment matters.\\
11 & foundation & 4.2.2 & Architecture is abstract; AD is a work product; method/notation/media are open.\\
12 & foundation & 4.2.4--4.2.5 & Viewpoint governs view; models are governed by model kinds.\\
13 & comparison & 4.2.6--4.2.7 & Correspondence and decision/rationale are first-class architecture-description concepts.\\
20--21 & comparison & 5.6--5.7 & Models can be shared across views; correspondences support consistency and traceability.\\
22 & comparison & 5.8 & Key decisions may record affected elements, alternatives, consequences, timestamps and citations.\\
24 & foundation & Clause 7 & Viewpoint conventions can include analysis methods and consistency operations.\\
28 & challenge & A.4 & Synthetic and projective view construction are both supported.\\
34 & foundation & B.2.6 & Model-kind metamodel lens: entities, attributes, relationships, constraints.\\
40 & comparison & C.4 & RM-ODP places information, interfaces, interactions and channel structure in different viewpoints.\\
\bottomrule\end{longtable}
\section{Chapter regression and status}
No Chapter 3 or Chapter 4 reopen trigger is found. The important downstream correction belongs to BA0/BA4: canonical-projective view generation is a candidate DDTA choice that must be justified, and `view` terminology should distinguish holistic concern-oriented views from partial models/renderings.
\begin{challenge}{Governance note}
This is a worksheet. The source PDF is copyrighted and is not included in the repository drop-in. The 2011 edition must not be used to claim that the 2022 edition contains identical normative clauses.
\end{challenge}
\end{document}
